\section{Реализация персонального ассистента организации времени обучающихся} 
\label{sec:development}

В данном разделе приведены обоснование используемых средств реализации, а также демонстрация её результатов. Для доказательства работоспособности системы и её соответствия требованиям был также разработан тестовый пример.

\subsection{Обоснование технических и программных средств разрабатываемой системы}

Выбор программных средств определяется архитектурными решениями, зафиксированными в разделе проектирования, а именно: модульным монолитом с трёхуровневым разделением ответственности, транзакционной изоляцией черновиков расписания и агентным модулем, совмещающим вызов языковой модели с детерминированным алгоритмом планирования. Каждое средство разработки соответствует конкретному слою системы и обосновывается техническими требованиями к нему.

\numberedparagraph
Серверная часть системы реализована на языке программирования Java версии 21, являющемся актуальным выпуском с долгосрочной поддержкой. Строгая статическая типизация языка обеспечивает надёжность доменной модели на этапе компиляции и снижает вероятность ошибок при приведении ответов языковой модели, поступающих в формате JSON, к объектам передачи данных. Зрелость экосистемы Java и наличие специализированных библиотек для интеграции с языковыми моделями также является определяющим фактором выбора.

В качестве основного каркаса приложения применяется Spring Boot версии 3.4. Данная платформа предоставляет механизмы автоконфигурации, встроенный веб-сервер и производственно-готовую инфраструктуру для построения REST-приложений без необходимости в ручной настройке большинства компонентов. Модуль Spring Web реализует слой контроллеров, обеспечивая маршрутизацию HTTP-запросов, десериализацию входящих данных и первичную валидацию полей объектов передачи данных.

Для разграничения доступа к ресурсам и управления сессиями используется модуль Spring Security совместно с механизмом токенов в формате JWT. Клиент передаёт токен в заголовке каждого запроса, а фильтр безопасности проверяет его подпись и срок действия до передачи управления контроллеру. Такой подход соответствует требованию безопасности REST-архитектуры, поскольку состояние сессии не хранится на сервере.

\numberedparagraph
В качестве реляционной системы управления базами данных выбрана PostgreSQL версии 17. Ключевым аргументом в её пользу является нативная поддержка ограничений исключения, позволяющих объявить непосредственно на уровне схемы правило, запрещающее сохранение перекрывающихся временных диапазонов для зафиксированных слотов одного пользователя. Данный механизм обеспечивает физическую защиту от коллизий расписания без дублирования логики проверки в прикладном коде. 

\begin{samepage}
\begin{lstlisting}[language=SQL, basicstyle=\footnotesize\ttfamily, frame=single,
    caption={SQL-запрос создания ограничения исключения для временных слотов},
    label={lst:overlap_constraint}]
CREATE EXTENSION IF NOT EXISTS btree_gist;

ALTER TABLE timeslots
ADD CONSTRAINT no_overlap
EXCLUDE USING gist (
    user_id WITH =,
    tsrange(start_time, end_time) WITH &&
) WHERE (is_committed = true);
\end{lstlisting}
\end{samepage}

Как видно из рисунка \ref{lst:overlap_constraint}, использование типа данных ``tsrange'' совместно с индексом ``GiST'' позволяет на уровне ядра базы данных физически запретить сохранение перекрывающихся временных интервалов для подтвержденных слотов одного пользователя.

Помимо этого, PostgreSQL предоставляет полноценную поддержку транзакций с настраиваемыми уровнями изоляции, необходимую для атомарного применения черновика расписания к основной таблице слотов.

Взаимодействие серверного приложения с базой данных организовано через модуль Spring Data JPA, реализующий паттерн «Репозиторий» посредством объектно-реляционного отображения. Применение данного модуля позволяет бизнес-логике оперировать исключительно объектами доменной модели, оставаясь независимой от конкретных запросов к базе данных, что соответствует архитектурному принципу слабой связности слоёв системы.

Управление версиями схемы базы данных производится инструментом Liquibase. При каждом запуске приложения инструмент последовательно применяет сценарии миграции, не применявшиеся ранее, что обеспечивает воспроизводимость структуры базы данных в любой среде развёртывания и исключает рассинхронизацию схемы с кодом приложения.

\numberedparagraph
Интеграция с внешним программным интерфейсом языковой модели реализована посредством библиотеки Spring AI версии 1.0. В качестве конкретного поставщика языковой модели используется сервис OpenRouter, предоставляющий унифицированный программный интерфейс к широкому спектру языковых моделей различных производителей. Совместимость OpenRouter с протоколом OpenAI API позволяет подключить его к Spring AI без изменений кода сервисного слоя --- достаточно переопределить базовый адрес запросов и ключ авторизации в конфигурационном файле приложения.

\begin{samepage}
\begin{lstlisting}[basicstyle=\footnotesize\ttfamily, frame=single,
    caption={Фрагмент конфигурации приложения для подключения OpenRouter},
    label={lst:spring_ai_config}]
spring.ai.openai.base-url=https://openrouter.ai/api
spring.ai.openai.api-key=${OPEN_ROUTER_API_KEY}
spring.ai.openai.chat.options.model=deepseek/deepseek-v4-flash

planning.ai.timeout-ms=15000
planning.ai.max-attempts=3
planning.ai.retry-backoff-ms=500
\end{lstlisting}
\end{samepage}

На рисунке \ref{lst:spring_ai_config} продемонстрировано, что благодаря совместимости протоколов, маршрутизация запросов к альтернативным языковым моделям (в данном случае к семейству DeepSeek) достигается исключительно конфигурационными механизмами. Дополнительно задаются параметры устойчивости (таймауты и количество попыток), необходимые для компенсации возможных сетевых задержек при обращении к внешнему API.

Обязательным требованием к выбираемой модели является поддержка режима структурированного вывода, поскольку именно он используется в реализации агентного модуля.

Режим структурированного вывода принудительно ограничивает генерацию модели заданной схемой данных, гарантируя синтаксическую корректность ответа. В реализации это используется для получения от языковой модели JSON-объекта, содержащего декомпозированные подзадачи и ограничения планирования. После получения ответа сервер выделяет JSON-фрагмент, выполняет его десериализацию в объект передачи данных и в случае ошибки разбора либо повторяет запрос, либо применяет упрощенный резервный сценарий.

На уровне прикладной логики агентный модуль получает от сервера текст запроса пользователя, календарный контекст, семантически релевантные фрагменты из векторного хранилища и историю текущего диалога. На основе этих данных языковая модель возвращает структурированный результат декомпозиции, который затем передается в детерминированный алгоритм распределения подзадач по временным окнам и далее в сервис жизненного цикла черновика.

\numberedparagraph
Выбор конкретной большой языковой модели для реализации агентного модуля основывался на сравнительном анализе легковесных моделей класса ``flash'', доступных через провайдера OpenRouter. Поскольку сложная задача удовлетворения ограничений (поиск свободных окон, расчет минут и проверка наложения дат) делегирована детерминированному алгоритму на стороне сервера (нейросимволический подход), необходимость в использовании дорогостоящих моделей с глубоким логическим выводом (reasoning) отсутствовала. Основными критериями выбора выступали: размер контекстного окна для интеграции подсистемы RAG, скорость ответа (latency), нативная поддержка структурированного вывода (Structured Outputs) и экономическая доступность. Сводные результаты сравнения моделей-кандидатов представлены в таблице~\ref{tab:llm_comparison}.

\begin{table}[!h!t]
\caption{Сравнительный анализ языковых моделей для агентного модуля}
\label{tab:llm_comparison}
\centering
\begin{tabularx}{\textwidth}{| >{\raggedright\arraybackslash}X | >{\centering\arraybackslash}p{0.22\textwidth} | >{\centering\arraybackslash}p{0.22\textwidth} | >{\centering\arraybackslash}p{0.22\textwidth} |}
\hline
\multicolumn{1}{|c|}{Критерий сравнения} & gpt-4o-mini & Claude 3.5 Haiku & DeepSeek-V3 \\ \hline
Размер контекстного окна & 128 000 токенов & 200 000 токенов & 128 000 токенов \\ \hline
Средняя задержка (latency) & Низкая (~0,56 с) & Очень низкая (~0,71 с) & Средняя (~0,60 с) \\ \hline
Ошибки синтаксиса JSON & Минимальные (~0,19 \%) & Низкие (~0,25 \%) & Допустимые (~0,51 \%) \\ \hline
Стоимость (за 1 млн входных токенов) & \$0,15 & \$0,80 & \$0,14 \\ \hline
\end{tabularx}
\end{table}

Как видно из таблицы~\ref{tab:llm_comparison}, модель ``Claude 3.5 Haiku'' предоставляет наибольший объем контекста, однако отличается высокой стоимостью эксплуатации. Модель ``gpt-4o-mini'' демонстрирует минимальный уровень синтаксических ошибок, но является закрытым проприетарным решением. В качестве базового вычислительного ядра была определена модель ``DeepSeek-V3'' (в конфигурации ``deepseek-chat''). Данная модель обеспечивает оптимальный баланс: достаточный объем контекстного окна для передачи расписания и методических материалов, низкую стоимость обработки запросов и открытую архитектуру весов. Когнитивного потенциала модели достаточно для качественного смыслового анализа и декомпозиции учебных задач, а интеграция через OpenRouter позволяет динамически переключаться на бесплатные демонстрационные эндпоинты на этапе тестирования без изменения программного кода системы.

\numberedparagraph
Реализация подсистемы семантического поиска в рамках архитектуры RAG потребовала разделения вычислительных задач на модель генерации (ответственную за структурированный вывод JSON) и специализированную модель векторизации (генерации эмбеддингов). Языковая модель генерации не способна эффективно выполнять поиск по неструктурированным массивам документов, в связи с чем текстовые фрагменты методических указаний предварительно преобразуются в числовые векторы признаков. Для этого используется легковесная модель эмбеддингов (например, ``text-embedding-3-small'' или локальная ``all-MiniLM-L6-v2''), кодирующая семантический смысл текстового чанка в многомерный массив чисел с плавающей точкой.

Для персистентного хранения и эффективной обработки полученных векторов применяется расширение ``pgvector'' для СУБД PostgreSQL. Данное решение позволяет сохранять эмбеддинги документов непосредственно в реляционной таблице в колонку типа ``vector'', исключая необходимость развертывания и поддержки обособленной векторной базы данных. При вызове метода поиска по сходству (``similaritySearch'') интерфейса ``VectorStore'' текстовый запрос пользователя векторизуется через ту же модель эмбеддингов, после чего расширение ``pgvector'' на уровне базы данных вычисляет математическое расстояние между вектором запроса и сохраненными фрагментами, используя метрику косинусного сходства (Cosine Similarity). Алгоритм извлекает строго заданное количество ``topK'' наиболее релевантных текстовых блоков, которые передаются в системный промпт в качестве изолированного семантического контекста, гарантируя точность декомпозиции и минимизируя риск галлюцинаций планирования.

\numberedparagraph
Пользовательский интерфейс реализован в виде одностраничного приложения на основе библиотеки React версии 18. Компонентная модель React обеспечивает декларативное управление состоянием интерфейса и позволяет реализовать интерактивные карточки черновика расписания как самостоятельные переиспользуемые компоненты, обновляющиеся при получении данных от сервера.

Для отображения календарной сетки применяется библиотека FullCalendar в редакции для React. Данная библиотека предоставляет готовый компонент с поддержкой представлений по неделям и дням, программного управления событиями и их перетаскивания в пределах сетки. Возможность перетаскивания событий непосредственно обеспечивает описанный в разделе проектирования сценарий ручного редактирования черновика расписания без дополнительного взаимодействия с агентом. Использование данной библиотеки исключает необходимость самостоятельной разработки логики позиционирования и отрисовки временных блоков.

Взаимодействие клиентского приложения с сервером организовано через библиотеку Axios. Перехватчики запросов данной библиотеки обеспечивают автоматическое добавление JWT-токена в заголовок каждого исходящего запроса, устраняя необходимость явного управления авторизацией в каждом компоненте интерфейса.

\numberedparagraph
Для обеспечения воспроизводимости среды разработки и упрощения локального развёртывания применяется платформа контейнеризации Docker совместно с инструментом Docker Compose. Конфигурационный файл Docker Compose описывает два сервиса: контейнер с серверным приложением и контейнер СУБД PostgreSQL. Такая конфигурация позволяет развернуть полностью работоспособное окружение без ручной установки зависимостей и настройки базы данных, что особенно важно для воспроизведения тестового примера при демонстрации работоспособности системы.

\subsection{Обзор результатов реализации}

В рамках реализации созданы серверное приложение на платформе Spring Boot, включающее слои контроллеров, бизнес-логики и доступа к данным, детерминированный алгоритм планирования временных слотов, механизм агентного взаимодействия с языковой моделью и клиентское одностраничное приложение с гибридным интерфейсом. В данном подразделе приведено описание ключевых нетривиальных решений, реализованных в ходе реализации.


\numberedparagraph

\numberedparagraph
Алгоритм автоматического распределения подзадач по временным интервалам реализует детерминированный жадный подход с пошаговым анализом календарной сетки. На вход алгоритм принимает упорядоченный список подзадач, массив уже зафиксированных слотов пользователя и объект ограничений, извлеченный языковой моделью (предпочтительное время суток, исключение выходных дней, дневной лимит нагрузки). Смещение расчетной точки во времени происходит итеративно с учетом пересечений с существующими календарными событиями. Псевдокод алгоритма приведен на рисунке~\ref{fig:algorithm_v2}.

\begin{samepage}
\begin{lstlisting}[basicstyle=\footnotesize\ttfamily, frame=single,
    caption={Псевдокод алгоритма планирования временных слотов с учетом ограничений},
    label={fig:algorithm_v2}]
Input: subtasks -- ordered list of subtasks
       occupied -- list of busy timeslots
       constraints -- {preferredTime, avoidWeekends, maxDailyMinutes}
Output: scheduledSlots

scheduledSlots <- []
cursor <- current_time
dailyBudgetUsed <- 0
currentDay <- cursor.date

for each subtask s in subtasks:
  placed <- false
  
  while cursor < deadline:
    cursor <- SnapToWorkingHours(cursor, constraints.preferredTime)
    
    if constraints.avoidWeekends AND IsWeekend(cursor):
      cursor <- NextWorkingDay(cursor)
      continue
      
    if (dailyBudgetUsed + s.duration) > constraints.maxDailyMinutes:
      cursor <- NextWorkingDay(cursor)
      dailyBudgetUsed <- 0
      continue
      
    conflict <- FindConflict(cursor, cursor + s.duration, occupied)
    if conflict != null:
      cursor <- conflict.end + MIN_GAP
      continue
      
    append ProposedTimeslot(s.id, cursor, cursor + s.duration) to scheduledSlots
    cursor <- cursor + s.duration + MIN_GAP
    dailyBudgetUsed <- dailyBudgetUsed + s.duration
    placed <- true
    break
\end{lstlisting}
\end{samepage}

Выбор временного интервала осуществляется без возврата: ни одно ранее принятое решение не пересматривается, что обеспечивает вычислительную сложность уровня $O(N \times M)$, где $N$ --- количество подзадач, а $M$ --- количество занятых слотов. Если свободного времени до установленного срока оказывается недостаточно, алгоритм прерывает цикл для текущего элемента, оставляя его без назначенного интервала. Такие подзадачи возвращаются пользователю в составе черновика для последующего ручного распределения в интерфейсе.

\numberedparagraph
Для обеспечения гибкого взаимодействия на естественном языке в агентном модуле реализована обработка двух функциональных сценариев: свободного текстового планирования и контекстно-зависимой декомпозиции. Сервисный слой динамически определяет логику инициализации объектов на основе наличия идентификатора родительской задачи в поступающем запросе. Ключевой фрагмент оркестрации данных сценариев в классе ``AgentServiceImpl'' представлен на рисунке~\ref{lst:agent_task_flow}.

\begin{samepage}
\begin{lstlisting}[language=Java, basicstyle=\footnotesize\ttfamily, frame=single,
    caption={Фрагмент оркестрации сценариев создания и контекстного связывания задач},
    label={lst:agent_task_flow}]
@Override
public DraftResponse processRequest(String message, UUID userId, UUID taskId) {
    Task task = taskId != null
            ? taskRepository.findById(taskId)
            .orElseThrow(() -> new EntityNotFoundException("Task not found: " + taskId))
            : null;

    String scheduleContext = buildScheduleContext(userId, task);
    String conversationId = conversationId(userId, taskId);
    reindexSemanticContext(userId, task);
    String semanticContext = buildSemanticContext(userId, task, message);
    String memoryContext = buildMemoryContext(conversationId);
    String currentDatetime = Instant.now().atOffset(ZoneOffset.UTC)
            .format(DateTimeFormatter.ISO_OFFSET_DATE_TIME);

    LlmPlannerResponse plannerResponse = invokeLlmWithRetry(
            message, taskId, scheduleContext, semanticContext, memoryContext, currentDatetime, task
    );

    LlmConstraints llmConstraints = plannerResponse.constraints();
    if (task == null) {
        task = createTaskFromPlannerResponse(userId, message, plannerResponse);
    }

    List<SubtaskDraft> drafts = toDrafts(plannerResponse.subtasks(), task.getId());
    AllocationConstraints constraints = toAllocationConstraints(llmConstraints, task);

    List<Timeslot> occupied = timeslotRepository
            .findByUserIdAndStartTimeGreaterThanEqualAndEndTimeLessThanEqual(
                    userId, constraints.rangeStart(), constraints.rangeEnd());

    List<ProposedTimeslot> proposedSlots = allocateSlotsSafely(drafts, occupied, constraints);
    drafts = mergeTimeslots(drafts, proposedSlots);

    chatMemory.add(conversationId, List.of(
            new UserMessage(message),
            new AssistantMessage(renderAssistantMemory(plannerResponse))
    ));

    return draftService.createDraftBatch(userId, drafts);
}
\end{lstlisting}
\end{samepage}

В рамках первого сценария (произвольный запрос пользователя) параметр ``taskId'' принимает значение ``null''. Сервис делегирует обработку сообщения языковой модели, которая на основе семантического анализа формулирует название новой глобальной цели в поле ``taskTitle'' и извлекает временной маркер в поле ``deadlineSemantic''. На основе полученной константы символьный слой серверной части (метод ``calculateDeadline'') детерминированно вычисляет точную дату окончания работы и выполняет персистентное сохранение сущности ``Task'' в реляционной базе данных. Все последующие сгенерированные подзадачи черновика автоматически связываются с идентификатором созданной на лету задачи.

Второй сценарий описывает декомпозицию в рамках уже существующего контекста. В данном случае клиентское приложение явно передает уникальный идентификатор ``taskId'', извлеченный из текущего экрана календарной сетки или внешней системы управления обучением. Сервис выполняет упреждающую выборку сущности из репозитория, извлекает жестко зафиксированные дедлайны и описание задачи, после чего передает их в системный промпт в качестве контекстного окружения. Языковая модель игнорирует построение новой родительской структуры, возвращая в поле ``taskTitle'' значение ``null'', и генерирует исключительно массив подзадач, который бесшовно привязывается к существующей цели без дублирования записей в базе данных. Подобное разделение ответственности позволяет поддерживать как концепцию гибридного управляемого интерфейса, так и свободный диалог с когнитивным агентом.

\numberedparagraph
Агентный модуль реализован в сервисе агента. Он принимает запрос пользователя, собирает структурированный календарный контекст из реляционной базы данных, извлекает дополнительный семантический контекст из векторного хранилища при наличии релевантных учебных материалов, обращается к языковой модели и затем преобразует полученный структурированный ответ в промежуточное представление подзадач, пригодное для дальнейшего алгоритмического планирования.

Календарный контекст включает сведения о текущей учебной задаче, дедлайнах и занятых временных интервалах пользователя. Эти данные передаются в системный промпт напрямую, поскольку имеют строгую структуру и требуют актуального состояния. Семантический контекст формируется отдельно на основе поиска по векторному хранилищу и используется только как вспомогательный источник информации о требованиях учебных материалов.

Для получения структурированного результата декомпозиции используется строго заданная схема JSON-ответа как на рисунке \ref{structured_out}.

\begin{samepage}
\begin{lstlisting}[language=Java, basicstyle=\footnotesize\ttfamily, frame=single,
    caption={Пример структурированного ответа модели},
    label={lst:structured_out}]
{
  "taskTitle": "Лабораторная работа по базам данных",
  "subtasks": [
    {
      "title": "Изучить методические указания и теорию по БД",
      "estimatedMinutes": 35,
      "cognitiveLoad": "MEDIUM"
    },
    {
      "title": "Спроектировать ER-диаграмму предметной области",
      "estimatedMinutes": 60,
      "cognitiveLoad": "HIGH"
    },
    {
      "title": "Написать SQL-скрипты DDL для создания таблиц",
      "estimatedMinutes": 45,
      "cognitiveLoad": "HIGH"
    },
    {
      "title": "Составить DML-запросы для заполнения таблиц данными",
      "estimatedMinutes": 40,
      "cognitiveLoad": "LOW"
    },
    {
      "title": "Оформить отчет со скриншотами выполнения",
      "estimatedMinutes": 45,
      "cognitiveLoad": "MEDIUM"
    }
  ],
  "constraints": {
    "deadlineSemantic": "END_OF_WEEK",
    "deadlineRawText": "к этим выходным",
    "priority": "HIGH",
    "preferredTimeOfDay": "ANY",
    "avoidWeekends": false,
    "maxDailyMinutes": null
  }
}
\end{lstlisting}
\end{samepage}

После возврата результата сервис выделяет JSON из текстового ответа модели, выполняет десериализацию в объект передачи данных и на его основе формирует набор подзадач и ограничений для последующего распределения по календарной сетке.

Пример ключевого фрагмента обработки структурированного вывода приведён на рисунке~\ref{lst:tool}.

\begin{samepage}
\begin{lstlisting}[language=Java, basicstyle=\footnotesize\ttfamily, frame=single,
    caption={Фрагмент обработки структурированного ответа языковой модели},
    label={lst:tool}]
LlmPlannerResponse plannerResponse = invokeLlmWithRetry(
        message, taskId, scheduleContext, semanticContext, currentDatetime, task
);

List<SubtaskDraft> drafts = toDrafts(plannerResponse.subtasks(), taskId);
AllocationConstraints constraints = toAllocationConstraints(plannerResponse.constraints(), task);
List<Timeslot> occupied = timeslotRepository
        .findByUserIdAndStartTimeGreaterThanEqualAndEndTimeLessThanEqual(
                userId, constraints.rangeStart(), constraints.rangeEnd());
List<ProposedTimeslot> proposedSlots = allocateSlotsSafely(drafts, occupied, constraints);

drafts = mergeTimeslots(drafts, proposedSlots);
return draftService.createDraftBatch(userId, drafts);
\end{lstlisting}
\end{samepage}

\numberedparagraph
Реализация подсистемы семантического поиска и архитектуры RAG в рамках серверной части приложения базируется на компонентах экосистемы Spring AI, предоставляющей декларативные абстракции для интеграции программного комплекса с технологиями искусственного интеллекта. Основным контрактом слоя данных для работы с эмбеддингами выступает интерфейс ``VectorStore''. В качестве конкретной реализации для СУБД PostgreSQL применяется класс ``PgVectorStore'', поставляемый специализированным стартером фреймворка. Интеграция с математическими моделями векторизации обеспечивается через интерфейс ``EmbeddingModel'', экземпляр которого внедряется в качестве обязательной зависимости в конструктор векторного хранилища. При выполнении операций добавления или поиска объектов класса ``Document'', компонент ``PgVectorStore'' автоматически делегирует текстовые фрагменты модели ``EmbeddingModel'' для генерации числовых векторов, после чего выполняет нативные SQL-запросы к базе данных. Пример программной конфигурации векторного хранилища представлен на рисунке~\ref{lst:vector_store_config}.

\begin{samepage}
\begin{lstlisting}[language=Java, basicstyle=\footnotesize\ttfamily, frame=single,
    caption={Программная конфигурация компонента PgVectorStore},
    label={lst:vector_store_config}]
@Configuration
public class VectorStoreConfiguration {

    @Bean
    public VectorStore vectorStore(JdbcTemplate jdbcTemplate, EmbeddingModel embeddingModel) {
        return PgVectorStore.builder()
                .jdbcTemplate(jdbcTemplate)
                .embeddingModel(embeddingModel)
                .dimensions(1536) 
                .distanceType(PgDistanceType.COSINE)
                .tableName("documents_vector")
                .build();
    }
}
\end{lstlisting}
\end{samepage}

Как показано в листинге~\ref{lst:vector_store_config}, конфигурирование векторного хранилища осуществляется с помощью шаблона проектирования ``Строитель'' (Builder). 
В процессе сборки компоненту передаются объект ``JdbcTemplate'' для выполнения низкоуровневых операций с реляционной базой данных и экземпляр ``EmbeddingModel''. 
Явная конфигурация позволяет жестко зафиксировать размерность генерируемых векторов (параметр ``dimensions'') и метрику определения семантического подобия (параметр ``distanceType''), 
в качестве которой выбрано косинусное расстояние.
% вставить формулу косинусного расстояния
Логический поток выполнения метода поиска по сходству полностью скрывает инфраструктурную сложность от сервисного слоя: переданная строка запроса преобразуется через ``EmbeddingModel'' в вектор, после чего ``PgVectorStore'' формирует SQL-запрос с использованием специализированных операторов расширения ``pgvector'' для извлечения заданного числа ``topK'' наиболее релевантных записей.

\numberedparagraph
Важным инфраструктурным компонентом модуля обработки документов является конвейер извлечения,
преобразования и загрузки данных, реализующий подготовку неструктурированных учебных материалов к последующему семантическому поиску.
Поступающий от клиентского приложения бинарный файл в составе объекта ``DocumentUploadRequest'' передается в сервисный слой, где трансформируется в набор изолированных векторизованных фрагментов.
Программная реализация данного процесса сосредоточена в методе ``uploadDocument'' класса ``DocumentServiceImpl'', ключевой фрагмент которого представлен на рисунке~\ref{lst:document_upload_impl}.

\begin{samepage}
\begin{lstlisting}[language=Java, basicstyle=\footnotesize\ttfamily, frame=single,
    caption={Реализация конвейера обработки и чанкинга документов в классе DocumentServiceImpl},
    label={lst:document_upload_impl}]
@Override
@Transactional
public DocumentResponse uploadDocument(UUID userId, DocumentUploadRequest req) {
    Resource resource = new ByteArrayResource(req.getContent());
    
    TikaDocumentReader documentReader = new TikaDocumentReader(resource);
    List<Document> rawDocuments = documentReader.get();

    for (Document doc : rawDocuments) {
        doc.getMetadata().put("userId", userId.toString());
        doc.getMetadata().put("filename", req.getFilename());
    }

    TokenTextSplitter textSplitter = new TokenTextSplitter(500, 350, 50, 10000, true);
    List<Document> chunkedDocuments = textSplitter.apply(rawDocuments);

    this.vectorStore.add(chunkedDocuments);

    return new DocumentResponse(
            UUID.randomUUID().toString(),
            req.getFilename(),
            req.getContentType(),
            Timestamp.from(Instant.now())
    );
}
\end{lstlisting}
\end{samepage}

Как показано в листинге на рисунке~\ref{lst:document_upload_impl}, на первом этапе бинарный контент интерпретируется с помощью компонента ``TikaDocumentReader'', обеспечивающего универсальное извлечение текстового содержимого из различных файловых форматов. Перед этапом сегментации документы обогащаются метаданными, содержащими идентификатор пользователя ``userId'', что критически важно для обеспечения изолированности учебного контекста в многопользовательской среде. На этапе преобразования (Transform) применяется токенизированный разделитель ``TokenTextSplitter''. Прямая подача полных текстов документов в векторное хранилище нецелесообразна из-за ограничений контекстного окна языковых моделей и снижения точности поиска. Алгоритм осуществляет разбиение текста на фрагменты (чанки) фиксированного размера (параметр ``chunkSize'', равный 500 токенам) с заданным перекрытием (параметр ``overlap'', равный 50 токенам). Наличие перекрытия предотвращает потерю семантических связей на границах разделения текстовых блоков. На финальном этапе загрузки (Load) сформированный массив чанков передается методу ``add'' интерфейса ``VectorStore'', который автоматически инициирует вызов модели эмбеддингов и персистентно сохраняет данные в таблицы СУБД PostgreSQL.

\numberedparagraph
Обеспечение изолированности предлагаемых изменений и реализация паттерна подтвержденного действия (Human-in-the-Loop)
потребовали разработки механизма транзакционного управления жизненным циклом черновика.
В процессе генерации расписания агентный модуль не модифицирует основную сетку календаря, а создает временные сущности: новые подзадачи получают статус ``DRAFT'', а предлагаемые временные интервалы помечаются флагом ``isCommitted = false''. 
В случае модификации уже существующей задачи, оригинальная сущность переводится в статус ``PENDING\_DELETE'', сохраняя свои данные для возможного отката. 

Ключевая логика фиксации и отмены изменений инкапсулирована в сервисе ``DraftServiceImpl''. На рисунке~\ref{lst:draft_lifecycle} представлены реализации методов подтверждения и отклонения черновика, демонстрирующие работу со статусной моделью.

\begin{samepage}
\begin{lstlisting}[language=Java, basicstyle=\footnotesize\ttfamily, frame=single,
    caption={Транзакционное управление состояниями черновика расписания},
    label={lst:draft_lifecycle}]
@Override
@Transactional
public List<SubtaskResponse> commitDraft(UUID userId) {
    List<Subtask> draftSubtasks = subtaskRepository.findByUserIdAndStatusIn(userId, DRAFT_STATUSES);
    List<Subtask> toSave = new ArrayList<>();
    List<Timeslot> timeslotsToCommit = new ArrayList<>();

    for (Subtask subtask : draftSubtasks) {
        if (subtask.getStatus() == SubtaskStatus.DRAFT) {
            subtask.setStatus(SubtaskStatus.SCHEDULED);
            toSave.add(subtask);
            for (Timeslot ts : subtask.getTimeslots()) {
                ts.setIsCommitted(true);
                timeslotsToCommit.add(ts);
            }
        } else if (subtask.getStatus() == SubtaskStatus.PENDING_DELETE) {
            subtaskRepository.delete(subtask);
        }
    }
    subtaskRepository.saveAll(toSave);
    timeslotRepository.saveAll(timeslotsToCommit);
    return toSave.stream().map(subtaskMapper::toResponse).toList();
}

@Override
@Transactional
public void rejectDraft(UUID userId) {
    subtaskRepository.deleteAllByUserIdAndStatus(userId, SubtaskStatus.DRAFT);
    List<Subtask> pendingDelete = subtaskRepository.findByUserIdAndStatus(userId, SubtaskStatus.PENDING_DELETE);
    
    for (Subtask subtask : pendingDelete) {
        subtask.setStatus(SubtaskStatus.SCHEDULED); 
    }
    subtaskRepository.saveAll(pendingDelete);
    timeslotRepository.deleteByUserIdAndIsCommitted(userId, false);
}
\end{lstlisting}
\end{samepage}

Как видно из листинга~\ref{lst:draft_lifecycle}, целостность данных гарантируется декларативным управлением транзакциями (аннотация ``@Transactional''). При вызове метода ``commitDraft'' система атомарно переводит все черновые подзадачи в статус запланированных (``SCHEDULED''), подтверждает временные слоты и физически удаляет элементы, ожидавшие замены. В случае вызова метода ``rejectDraft'', происходит обратный процесс (откат): временные черновые сущности удаляются из базы данных, а исходные подзадачи, помеченные к удалению, восстанавливаются в статус ``SCHEDULED''.
Подобная реализация конечного автомата состояний в сочетании с транзакционным контекстом базы данных полностью исключает риск частичного применения расписания или потери данных пользователя в случае программных сбоев.

\subsection{Демонстрация работоспособности системы}
% Для демонстрации работоспособности разработанной системы был подготовлен тестовый сценарий, охватывающий полный жизненный цикл взаимодействия пользователя с интеллектуальным ассистентом: исходное состояние интерфейса, генерацию черновика расписания, его ручное редактирование, подтверждение, а также обработку частично невыполнимого плана и отклонение черновика. Все изображения получены на основе прототипа пользовательского интерфейса, воспроизводящего спроектированные функции системы.

% Начальное состояние интерфейса представлено на рисунке \ref{fig:screen1}. Экран содержит диалоговую панель ассистента, недельную календарную сетку с уже занятыми интервалами и информационный блок со сводкой состояния. На данном этапе черновик ещё не сформирован, однако пользователь уже может инициировать сценарий планирования посредством текстового запроса.

%фикс: карточка создания задачи на фоне деолтного скрина. создается лекция на 12.00-13.10
\begin{figure}[H]
\centering
\includegraphics[width=1.0\linewidth]{picture/screens/manual_task_creation.png}
\caption{Главный экран с карточкой создания подзадачи}
\label{fig:screen1}
\end{figure}

%фикс: стартовый скриншот с двумя задачами: лекции в понедельник (12.00-13.20 и 13.50-15.10)
\begin{figure}[H]
\centering
\includegraphics[width=1.0\linewidth]{picture/screens/default_start_screen.png}
\caption{Главный экран системы до генерации черновика расписания}
\label{fig:screen1}
\end{figure}

% После отправки запроса ассистент выполняет декомпозицию учебной задачи и инициирует алгоритм подбора свободных интервалов. Результат отображается в виде карточки черновика и набора временных блоков в календаре, выделенных отдельным цветом, что показано на рисунке \ref{fig:screen2}. Таким образом, пользователь получает возможность сопоставить текстовое описание предложенного плана с его визуальным представлением в календарной сетке.

%фикс: преложенный черновик расписания из пяти подзадач для лабораторонй по бд
\begin{figure}[H]
\centering
\includegraphics[width=1.0\linewidth]{picture/screens/propose_draft.png}
\caption{Формирование черновика расписания на основе запроса пользователя}
\label{fig:screen2}
\end{figure}

% Следующий этап демонстрирует возможность ручной корректировки автоматически построенного плана. Как видно из рисунка \ref{fig:screen3}, пользователь может изменить положение отдельного временного блока и исключить одну из подзадач из текущего черновика. Данный сценарий подтверждает, что разработанный интерфейс поддерживает паттерн Human-in-the-Loop и не навязывает пользователю автоматически сгенерированное решение без возможности его ручной доработки.

%пользователь перетягивает карточку с задачей написания er-диаграммы со вторника на среду
\begin{figure}[H]
\centering
\includegraphics[width=1.0\linewidth]{picture/screens/edit_draft_manually.png}
\caption{Пример редактирования черновика расписания пользователем}
\label{fig:screen3}
\end{figure}

% После подтверждения черновика его элементы фиксируются, а соответствующие временные интервалы включаются в основное расписание пользователя. На рисунке \ref{fig:screen4} показано изменение статуса задачи и визуального оформления временных слотов после подтверждения предложенного варианта.

%коммит черновика + отметки о сделнных заадчах + прогресс бар справа для геймификации и визуализации прогресса
\begin{figure}[H]
\centering
\includegraphics[width=1.0\linewidth]{picture/screens/commit_draft_and_completed_tasks_and_gamification.png}
\caption{Интерфейс системы при подтверждении черновика расписания пользователем. Дополнительно изображена шкала прогресса на основе выполненных задач}
\label{fig:screen4}
\end{figure}

%Отдельно был рассмотрен граничный сценарий, в котором свободных временных окон до установленного срока недостаточно для размещения всех подзадач. Как показано на рисунке \ref{fig:screen5}, система формирует частичный черновик и явно отображает список элементов, которые не удалось автоматически встроить в календарь. Такое поведение позволяет пользователю своевременно скорректировать дедлайн, вручную разместить проблемные подзадачи или отказаться от предложенного варианта.
%фикс: невозможность вставки при скором дедлайне
\begin{figure}[H]
\centering
\includegraphics[width=1.0\linewidth]{picture/screens/dedline_conflict.png}
\caption{Пример изображения уведомления о конфликте черновика при недостатке свободных интервалов}
\label{fig:screen5}
\end{figure}

%Дополнительно был рассмотрен сценарий отклонения ранее сформированного черновика. На рисунке \ref{fig:screen6} показано системное уведомление, подтверждающее отмену предложенного плана и возможность формирования нового варианта. Данный сценарий демонстрирует корректную работу механизма отката без фиксации временных слотов в основном расписании.
%фикс: отклонение черновика для плана эссе
\begin{figure}[H]
\centering
\includegraphics[width=1.0\linewidth]{picture/screens/reject_draft_screen.png}
\caption{Пример сообщения об отклонении черновика расписания}
\label{fig:screen6}
\end{figure}

\begin{figure}[H]
\centering
\includegraphics[width=1.0\linewidth]{picture/screens/rag_draft_economics_method.png}
\caption{Пример черновика на основе методического документа по экономическому обоснованию}
\label{fig:screen6}
\end{figure}

\begin{figure}[H]
\centering
\includegraphics[width=1.0\linewidth]{picture/screens/rag_draft_economics_method_specific_subtask.png}
\caption{Пример карточки созданной подзадачи, с подробной специфической информацией}
\label{fig:screen6}
\end{figure}

Представленные примеры отражают основные сценарии работы системы: формирование черновика, визуальное представление результата работы агентного модуля, ручную корректировку временных блоков, подтверждение изменений и их отклонение, а также информирование пользователя о невозможности полного планирования при текущих ограничениях календаря.


\subsection{Вывод}

В результате реализации создан программный комплекс персонального ассистента организации времени обучающихся, соответствующий архитектурным решениям и функциональным требованиям, зафиксированным в разделе проектирования.

Серверная часть системы построена в виде модульного монолита на платформе Spring Boot с трёхуровневым разделением ответственности. Слой контроллеров обеспечивает приём и первичную валидацию запросов, сервисный слой реализует бизнес-логику, включая управление жизненным циклом черновика и координацию вызовов агентного модуля, а слой доступа к данным инкапсулирует взаимодействие с СУБД через интерфейсы репозиториев.

Нейросимволический подход к планированию реализован посредством чёткого разделения ответственности между двумя компонентами. Языковая модель, интегрированная через Spring AI с поставщиком OpenRouter, выполняет декомпозицию задачи и возвращает структурированный объект передачи данных. Жадный алгоритм размещения принимает полученные подзадачи и распределяет их по свободным временным окнам с учётом рабочего диапазона дня, допустимых перерывов, дневных ограничений и занятости пользователя.

Транзакционная целостность данных обеспечена на двух уровнях: на уровне приложения --- за счёт изоляции черновых изменений и их атомарного подтверждения или отката; на уровне базы данных --- за счёт ограничений исключения PostgreSQL, физически исключающих пересечение зафиксированных временных диапазонов при подтверждении черновика.

Клиентское одностраничное приложение реализует гибридный агентный интерфейс: диалоговое окно позволяет формулировать намерения в произвольной текстовой форме, а интерактивная календарная сетка на основе FullCalendar обеспечивает визуальное редактирование и подтверждение сгенерированного черновика расписания. Воспроизводимость среды развёртывания обеспечена средствами Docker Compose.
